\def\probtitle{사전순 최소}
\def\probno{D} % 문제 번호

\begin{problem}{\probtitle}{표준 입력(stdin)}{표준 출력(stdout)}{1 초}{1024 megabytes}

길이 $N$인 문자열 $T$가 주어진다. $T$는 알파벳 소문자로만 이루어져 있다.

당신은 다음 연산을 $0$회 이상 사용할 수 있다.

\begin{itemize}[topsep=0pt,noitemsep]
    \item 두 정수 $l, r$를 선택한다. ($1 \le l < r \le |T|$)
    \item 만약 부분 문자열 $T_l T_{l+1} \dots T_r$의 모든 문자가 동일하다면, 그 구간에서 문자 하나만 남기고 나머지를 모두 제거한다.
\end{itemize}

즉, $T_l = T_{l+1} = \cdots = T_r = \scalebox{1.3}{c}$ 인 경우,

문자열
$ T_1 T_2 \dots T_{l-1}\, \underline{T_l T_{l+1} \dots T_r}\, T_{r+1} \dots T_{|T|} $
를
$ T_1 T_2 \dots T_{l-1}\, \underline{\scalebox{1.3}{c}}\, T_{r+1} \dots T_{|T|} $
로 바꿀 수 있다.

이 연산을 0회 이상 원하는 만큼 사용하여 만들 수 있는 모든 문자열 중, 사전 순으로 가장 앞서는 문자열을 구하여라.

\InputFile

첫째 줄에 문자열 $T$의 길이를 나타내는 정수 $N$이 주어진다. 

둘째 줄에 길이가 $N$인 문자열 $T$가 주어진다.


\OutputFile

가능한 모든 결과 문자열 중 사전 순으로 가장 앞서는 문자열을 출력한다.

\Constraints

\begin{itemize}[topsep=0pt,noitemsep]
    \item $1 \le N \le 200\,000$
    \item $T$는 알파벳 소문자로만 이루어져 있다.
\end{itemize}

\Example

\begin{example}
    \exmpfile{./example/01.in.txt}{./example/01.out.txt}%
    \exmpfile{./example/02.in.txt}{./example/02.out.txt}%
\end{example}

\newpage

\Notes

사전 순 비교는 다음과 같이 정의한다.

두 문자열 $A$, $B$에 대하여, 다음 둘 중 하나를 만족하면 $A$가 $B$보다 사전 순으로 앞선다고 한다.

\begin{itemize}[topsep=0pt,noitemsep]
    \item $A$가 $B$의 접두사인 경우
    \item 처음으로 다른 위치를 $i$라고 할 때, $A_i$가 $B_i$보다 알파벳 순으로 앞서는 경우
\end{itemize}

예를 들어, \texttt{ab}는 \texttt{aba}보다 사전 순으로 앞서고, \texttt{abc}는 \texttt{abd}보다 사전 순으로 앞선다.

어떤 구간에 연산을 적용할 수 있더라도 반드시 적용해야 하는 것은 아니다.

예제 1의 문자열 \texttt{aabbba}의 경우, 만들 수 있는 문자열은
\texttt{aabbba}, \texttt{abbba}, \texttt{aabba}, \texttt{abba}, \texttt{aaba}, \texttt{aba}
의 6가지이다.

이 중 사전 순으로 가장 앞서는 문자열은 \texttt{aaba}이다.

\end{problem}